\documentclass[12pt,oneside]{book}

\usepackage[a4paper,margin=25mm]{geometry}
\usepackage{array}
\usepackage{booktabs}
\usepackage{enumitem}
\usepackage{hyperref}
\usepackage{longtable}
\usepackage{tabularx}
\usepackage{xurl}

\hypersetup{
    colorlinks=true,
    linkcolor=blue,
    urlcolor=blue,
    citecolor=blue
}

\setlist[itemize]{leftmargin=*,topsep=4pt,itemsep=2pt}
\setlist[enumerate]{leftmargin=*,topsep=4pt,itemsep=2pt}

\title{\textbf{Underground Design and Construction}\\
Draft Textbook Content Outline for Review}
\author{Compiled from RSE3010 Weekly Study Notes and Lecture Slides}
\date{\today}

\begin{document}
\pagestyle{plain}

\maketitle
\tableofcontents

\chapter*{Review Basis and Proposed Logic}
\addcontentsline{toc}{chapter}{Review Basis and Proposed Logic}

\section*{Local Document Basis}
\addcontentsline{toc}{section}{Local Document Basis}

This outline is compiled from the two project folders:

\begin{itemize}
    \item \texttt{Weekly Study Notes}: Weeks 1--9 were used as the main detailed text basis.
    \item \texttt{Lecture Slides}: Weeks 1--9 and Week 12 were used to check lecture emphasis and add the construction-methods chapter.
    \item Week 11 material is intentionally excluded according to the previous review instruction.
\end{itemize}

\section*{Benchmarking Against Comparable Textbooks}
\addcontentsline{toc}{section}{Benchmarking Against Comparable Textbooks}

The proposed book structure is not a direct week-by-week copy. It follows a textbook logic informed by comparable published books:

\begin{itemize}
    \item Rock engineering texts typically build from geological setting, stress, intact rock, discontinuities, rock masses, classification, excavation principles and support principles before detailed underground design.
    \item Tunnel engineering handbooks usually move early into tunnel planning, geotechnical investigation, stabilisation, lining, excavation methods, TBMs, shotcrete, immersed tunnels, cut-and-cover structures, services, safety and contracting.
    \item Tunnelling mechanics books commonly separate practical design topics from mechanics topics such as ground behaviour, stress redistribution, face stability, support action and tunnel response.
    \item ADECO-RS literature emphasises a design-to-construction sequence: survey, diagnosis, therapy, operation, monitoring and final adjustment.
    \item Practical construction books emphasise method selection, drill-and-blast, TBM, NATM, NMT, roadheaders, ground control, mucking and haulage, grouting, water handling and lining.
\end{itemize}

\section*{Recommended Textbook Logic}
\addcontentsline{toc}{section}{Recommended Textbook Logic}

The recommended logic is:

\begin{enumerate}
    \item Explain why underground engineering matters and how project risk is framed.
    \item Establish the ground model and uncertainty before using parameters in design.
    \item Build the rock mechanics basis: stress, intact rock and failure criteria.
    \item Convert geological observations into rock mass classifications and design parameters.
    \item Identify likely excavation behaviour and select preliminary support.
    \item Check the design with analytical and numerical methods.
    \item Integrate design evidence into a coherent recommendation.
    \item Close the loop with construction method selection, support installation and monitoring.
\end{enumerate}

\section*{Source-to-Chapter Map}
\addcontentsline{toc}{section}{Source-to-Chapter Map}

\begin{longtable}{p{0.24\textwidth}p{0.68\textwidth}}
\toprule
\textbf{Primary source} & \textbf{Proposed textbook use} \\
\midrule
Week 1 notes and slides & Chapter 1: infrastructure context, project risk, shallow and deep underground conditions, environmental assessment, Australian case studies. \\
Week 5 notes and slides & Chapter 2: site investigation, RQD, borehole interpretation, uncertainty, modern investigation methods and data fusion. \\
Week 2 notes and slides & Chapter 3: stress, intact rock mechanics, Mohr--Coulomb, Hoek--Brown, tensile cut-off and equivalent parameter fitting. \\
Week 3 notes and slides & Chapter 4: RMR, Q-system, GSI, classification workflow and final GSI selection. \\
Week 4 notes and slides & Chapter 5: generalised Hoek--Brown, disturbance factor, rock mass strength, modulus and equivalent Mohr--Coulomb parameters. \\
Week 6 notes and slides & Chapter 6: failure mechanisms and empirical support design using RMR and Q-system. \\
Week 7 notes and slides & Chapter 7: Kirsch solutions, elastic-plastic response, convergence--confinement method, GRC, SCC, LDP and support timing. \\
Week 8 notes and slides & Chapter 8: settlement prediction, Mair method, FEM workflow, RS2 modelling, lining modelling and practical numerical checks. \\
Week 9 notes and slides & Chapter 9: integrated design workflow, uncertainty synthesis and method comparison. \\
Week 12 slides & Chapter 10: conventional, mechanised, surface and immersed construction; NATM, NMT, CCM, ADECO-RS; support systems and monitoring. \\
\bottomrule
\end{longtable}

\part{Project Context, Ground Models and Design Basis}

\chapter{Underground Engineering in Infrastructure and Energy Projects}

\section{Chapter Purpose}
Introduce underground design as a risk-based engineering process shaped by infrastructure demand, geological uncertainty, environmental approval, construction access and long-term asset performance.

\section{Core Content Outline}
\begin{enumerate}
    \item Infrastructure market capacity and the demand for geotechnical engineers.
    \item Types of underground works: transport tunnels, water tunnels, energy storage, station caverns and underground power facilities.
    \item Major Australian project context:
    \begin{itemize}
        \item Melbourne Metro Tunnel.
        \item Snowy 2.0.
        \item West Gate Tunnel.
        \item Sydney Metro and SRL East as recurring examples.
    \end{itemize}
    \item Infrastructure risk dashboards and the role of geotechnical uncertainty.
    \item Shallow versus deep underground conditions:
    \begin{itemize}
        \item Ground movement and serviceability for shallow urban works.
        \item In-situ stress, spalling, squeezing and access constraints for deep works.
    \end{itemize}
    \item Environmental assessment frameworks:
    \begin{itemize}
        \item EIS and EES.
        \item EPBC Act triggers and approval interfaces.
    \end{itemize}
    \item Net zero, pumped hydro and the renewed value of underground space.
    \item From project context to design requirements.
\end{enumerate}

\section{Proposed Figures and Tables}
\begin{itemize}
    \item Workflow figure: project drivers to ground model to design method to construction decision.
    \item Comparison table: shallow urban tunnel versus deep hard-rock tunnel.
    \item Risk dashboard table for geotechnical, environmental and construction risks.
\end{itemize}

\section{Worked Examples and Exercises}
\begin{itemize}
    \item Compare Melbourne Metro and Snowy 2.0 in terms of design drivers.
    \item Explain how poor ground knowledge affects cost, program and support decisions.
\end{itemize}

\chapter{Site Investigation, Ground Models and Uncertainty}

\section{Chapter Purpose}
Move site investigation earlier in the book so readers understand where rock mechanics inputs come from before using them in design equations.

\section{Core Content Outline}
\begin{enumerate}
    \item Role of site investigation in underground design.
    \item Ground model development:
    \begin{itemize}
        \item Geological units.
        \item Weathering zones.
        \item Structural domains.
        \item Groundwater conditions.
        \item Existing structures and construction sensitivity.
    \end{itemize}
    \item Borehole planning and core logging.
    \item Rock Quality Designation:
    \begin{itemize}
        \item Definition and calculation.
        \item Interpretation and limitations.
        \item Drill orientation bias.
        \item Statistical variability.
    \end{itemize}
    \item SRL East borehole interpretation.
    \item From investigation data to design parameters.
    \item Parameter uncertainty and sensitivity:
    \begin{itemize}
        \item GSI sensitivity in Hoek--Brown calculations.
        \item Scale effects.
        \item Data quality and confidence ratings.
    \end{itemize}
    \item Modern investigation methods:
    \begin{itemize}
        \item Digital core logging.
        \item LiDAR and photogrammetry in tunnels.
        \item Geophysics and data fusion.
        \item Monitoring feedback and digital ground models.
    \end{itemize}
\end{enumerate}

\section{Proposed Figures and Tables}
\begin{itemize}
    \item Ground model development flowchart.
    \item RQD calculation sketch and limitations table.
    \item Data confidence matrix linking investigation quality to design conservatism.
\end{itemize}

\section{Worked Examples and Exercises}
\begin{itemize}
    \item Calculate RQD from a core run and discuss uncertainty.
    \item Interpret simplified SRL East borehole logs for design implications.
\end{itemize}

\part{Rock Mechanics and Rock Mass Design Parameters}

\chapter{Stress, Intact Rock Mechanics and Failure Criteria}

\section{Chapter Purpose}
Provide the mechanical foundation for stress analysis, intact rock failure and the selection of strength criteria used later for tunnel design.

\section{Core Content Outline}
\begin{enumerate}
    \item Overview of material failure theory.
    \item Stress at a point and principal stresses.
    \item In-situ stress:
    \begin{itemize}
        \item Vertical stress from overburden.
        \item Horizontal stress and tectonic components.
        \item Elastic at-rest approximation.
        \item Stress regimes and design implications.
    \end{itemize}
    \item Selecting a failure criterion for underground design.
    \item Mohr--Coulomb failure criterion:
    \begin{itemize}
        \item Basic equation.
        \item Cohesion and friction angle.
        \item Uniaxial and tensile strengths.
        \item Failure plane orientation.
        \item Mohr circle interpretation.
        \item Limitations.
    \end{itemize}
    \item Hoek--Brown failure criterion for intact rock:
    \begin{itemize}
        \item Intact rock equation.
        \item Material constant \(m_i\).
        \item Tensile strength from Hoek--Brown.
        \item Worked example for marble.
    \end{itemize}
    \item Converting Hoek--Brown to equivalent Mohr--Coulomb parameters:
    \begin{itemize}
        \item Fitting range.
        \item Mohr circle regression.
        \item Linear regression in \(\sigma_1^f-\sigma_3\) space.
    \end{itemize}
    \item Griffith theory and tensile cut-off.
    \item Notation discipline: in-situ principal stresses versus failure stresses.
\end{enumerate}

\section{Proposed Figures and Tables}
\begin{itemize}
    \item In-situ stress components and principal stress notation.
    \item Mohr--Coulomb envelope and Hoek--Brown curve comparison.
    \item Parameter fitting workflow.
\end{itemize}

\section{Worked Examples and Exercises}
\begin{itemize}
    \item Calculate vertical and horizontal stress at depth.
    \item Estimate confined compressive strength using Hoek--Brown.
    \item Convert Hoek--Brown data into equivalent Mohr--Coulomb parameters.
\end{itemize}

\chapter{Rock Mass Classification: RMR, Q-System and GSI}

\section{Chapter Purpose}
Convert geological observations into reproducible engineering indices that support communication, parameter derivation and preliminary support design.

\section{Core Content Outline}
\begin{enumerate}
    \item Why rock mass classification is used.
    \item Classification workflow from logging to design input.
    \item Rock Mass Rating system:
    \begin{itemize}
        \item Intact rock strength.
        \item RQD.
        \item Discontinuity spacing.
        \item Discontinuity condition.
        \item Groundwater condition.
        \item Rock mass classes.
    \end{itemize}
    \item NGI Q-system:
    \begin{itemize}
        \item \(RQD/J_n\): block size.
        \item \(J_r/J_a\): interblock shear strength.
        \item \(J_w/SRF\): active stress and water effects.
        \item Rock mass quality classes.
    \end{itemize}
    \item Geological Strength Index:
    \begin{itemize}
        \item Structure descriptions.
        \item Joint surface condition.
        \item Chart-based GSI estimation.
        \item Empirical GSI correlations.
    \end{itemize}
    \item Selecting a final GSI value:
    \begin{itemize}
        \item Reconciliation of RMR, Q and chart estimates.
        \item Conservatism and uncertainty.
        \item Documentation for design review.
    \end{itemize}
\end{enumerate}

\section{Proposed Figures and Tables}
\begin{itemize}
    \item Classification workflow diagram.
    \item RMR, Q and GSI comparison table.
    \item GSI selection decision table.
\end{itemize}

\section{Worked Examples and Exercises}
\begin{itemize}
    \item Complete RMR and Q-system calculations for a rock mass domain.
    \item Compare multiple GSI estimates and justify a final design value.
\end{itemize}

\chapter{Rock Mass Strength, Deformability and Design Parameters}

\section{Chapter Purpose}
Translate rock mass quality into strength and deformability parameters suitable for empirical, analytical and numerical underground design.

\section{Core Content Outline}
\begin{enumerate}
    \item Position of design parameter derivation in the underground design chain.
    \item Generalised Hoek--Brown criterion for rock masses.
    \item Rock mass Hoek--Brown parameters:
    \begin{itemize}
        \item \(m_b\).
        \item \(s\).
        \item \(a\).
        \item \(\sigma_{ci}\), \(m_i\), GSI and disturbance factor.
    \end{itemize}
    \item Disturbance factor:
    \begin{itemize}
        \item TBM excavation.
        \item Controlled blasting.
        \item Poor drill-and-blast excavation.
        \item Construction method sensitivity.
    \end{itemize}
    \item Rock mass strength parameters:
    \begin{itemize}
        \item Uniaxial compressive strength of rock mass.
        \item Triaxial strength at in-situ confinement.
        \item Tensile strength of rock mass.
    \end{itemize}
    \item Deformation modulus of rock mass:
    \begin{itemize}
        \item Hoek--Diederichs generalised equation.
        \item Hoek--Diederichs simplified equation.
        \item Hoek et al. GSI--UCS equation.
        \item RMR-based equations.
        \item Q-based equation.
    \end{itemize}
    \item Equivalent Mohr--Coulomb parameters for rock masses.
    \item Parameter comparison and sensitivity tables.
\end{enumerate}

\section{Proposed Figures and Tables}
\begin{itemize}
    \item Hoek--Brown rock mass parameter workflow.
    \item Disturbance factor comparison table.
    \item Strength and modulus calculation templates.
\end{itemize}

\section{Worked Examples and Exercises}
\begin{itemize}
    \item Derive \(m_b\), \(s\), \(a\), \(\sigma_{cm}\), \(\sigma_t\) and \(E_m\) for a rock mass.
    \item Compare TBM and poor drill-and-blast disturbance assumptions.
    \item Convert rock mass Hoek--Brown to equivalent Mohr--Coulomb parameters.
\end{itemize}

\part{Underground Excavation Behaviour and Support Design}

\chapter{Underground Failure Mechanisms and Empirical Support Design}

\section{Chapter Purpose}
Identify the likely failure mechanism before selecting preliminary support using rock mass classification systems.

\section{Core Content Outline}
\begin{enumerate}
    \item Linking investigation uncertainty to support design.
    \item Main underground failure mechanisms:
    \begin{itemize}
        \item Stress-induced failure.
        \item Structurally controlled failure in underground openings.
        \item Large deformation and squeezing ground.
    \end{itemize}
    \item Failure-mechanism decision workflow.
    \item Empirical support design using RMR:
    \begin{itemize}
        \item Unsupported span.
        \item Stand-up time.
        \item Adjustment for joint orientation.
        \item Support recommendations.
    \end{itemize}
    \item Empirical support design using the Q-system:
    \begin{itemize}
        \item Excavation Support Ratio.
        \item Equivalent dimension \(D_e\).
        \item Q-chart workflow.
        \item Support categories.
        \item Wall support adjustment.
        \item Maximum unsupported span.
    \end{itemize}
    \item Case study: railway station cavern.
    \item Case note: Sydney Opera House car park.
    \item Influence of input uncertainty on support class.
\end{enumerate}

\section{Proposed Figures and Tables}
\begin{itemize}
    \item Failure mechanism decision tree.
    \item RMR support recommendation table.
    \item Simplified Q-system support category table.
    \item Support sensitivity table for uncertain RMR and Q inputs.
\end{itemize}

\section{Worked Examples and Exercises}
\begin{itemize}
    \item Estimate RMR-based stand-up time and preliminary support.
    \item Design Q-system support for tunnel and cavern spans.
    \item Discuss why two empirical systems may produce different support recommendations.
\end{itemize}

\chapter{Analytical Tunnel Design and Ground--Support Interaction}

\section{Chapter Purpose}
Introduce closed-form and semi-analytical methods that check stress redistribution, deformation, plastic zone development and support timing.

\section{Core Content Outline}
\begin{enumerate}
    \item Assumptions and limitations of analytical methods.
    \item Elastic stress solutions:
    \begin{itemize}
        \item Kirsch equations.
        \item Boundary stresses at the tunnel wall.
        \item Stress concentration around circular openings.
        \item Engineering applications and limitations.
    \end{itemize}
    \item Elastic-plastic response:
    \begin{itemize}
        \item Rationale for plasticity analysis.
        \item Critical internal pressure.
        \item Plastic zone radius for unsupported openings.
        \item Stress redistribution in the plastic zone.
    \end{itemize}
    \item Convergence--confinement method:
    \begin{itemize}
        \item Conceptual reduction of internal support pressure.
        \item Ground Reaction Curve.
        \item Support Characteristic Curve.
        \item Longitudinal Deformation Profile.
        \item Installation timing and \(u_{s0}\).
        \item Equilibrium point and factor of safety.
    \end{itemize}
    \item Support stiffness, support capacity and deformation allowance.
    \item Role of spreadsheet calculation and RocSupport-type tools.
    \item Case study: Gotthard Motorway Tunnel.
    \item Worked example: Snowy 2.0 at depth.
\end{enumerate}

\section{Proposed Figures and Tables}
\begin{itemize}
    \item Kirsch stress distribution around a circular tunnel.
    \item GRC--SCC--LDP combined design diagram.
    \item Support installation timing comparison table.
\end{itemize}

\section{Worked Examples and Exercises}
\begin{itemize}
    \item Calculate tunnel wall stresses using Kirsch equations.
    \item Estimate critical pressure and plastic zone radius.
    \item Construct GRC, SCC and LDP curves and find the equilibrium point.
\end{itemize}

\chapter{Numerical Modelling and Settlement Prediction}

\section{Chapter Purpose}
Show how empirical and analytical designs are checked against settlement, staged excavation, geometry, lining response and material behaviour using numerical models.

\section{Core Content Outline}
\begin{enumerate}
    \item Ground movement caused by tunnelling.
    \item Gaussian settlement trough:
    \begin{itemize}
        \item Peck-type settlement trough.
        \item Volume loss.
        \item Maximum settlement.
        \item Shallow and deep tunnel response.
    \end{itemize}
    \item Components of ground movement in TBM tunnelling.
    \item Material properties for assignment-style modelling.
    \item Tunnel geometry and geological idealisation.
    \item Initial stress worked example.
    \item Mair method worked examples.
    \item Finite element modelling principles:
    \begin{itemize}
        \item Continuum idealisation.
        \item Constitutive models.
        \item Boundary conditions.
        \item Mesh refinement.
        \item Failure criteria.
    \end{itemize}
    \item RS2 modelling workflow:
    \begin{itemize}
        \item Geometry.
        \item Materials.
        \item Staging.
        \item Lining model.
        \item Result extraction.
    \end{itemize}
    \item Practical modelling issues:
    \begin{itemize}
        \item Shallow tunnel convergence problems.
        \item Lining simplification.
        \item Groundwater.
        \item Model dimensions.
        \item Interpreting non-convergence.
    \end{itemize}
    \item Comparing Mair method and RS2 results.
    \item Digital transformation: BIM, digital twins and monitoring updates.
\end{enumerate}

\section{Proposed Figures and Tables}
\begin{itemize}
    \item Gaussian settlement trough schematic.
    \item FEM modelling workflow.
    \item Modelling assumptions checklist.
    \item Mair method versus RS2 comparison table.
\end{itemize}

\section{Worked Examples and Exercises}
\begin{itemize}
    \item Estimate surface settlement from tunnel depth, diameter and volume loss.
    \item Build a simplified FEM modelling plan.
    \item Compare numerical output with analytical predictions and explain differences.
\end{itemize}

\chapter{Integrated Underground Design}

\section{Chapter Purpose}
Bring the previous methods together into a coherent engineering recommendation, with explicit treatment of uncertainty and disagreement between methods.

\section{Core Content Outline}
\begin{enumerate}
    \item The integrated design workflow:
    \begin{itemize}
        \item Site investigation.
        \item Empirical design.
        \item Analytical design.
        \item Numerical verification.
        \item Construction and monitoring feedback.
    \end{itemize}
    \item How methods complement each other.
    \item What each method can and cannot answer.
    \item Resolving disagreement between empirical, analytical and numerical results.
    \item Uncertainty through the workflow:
    \begin{itemize}
        \item Ground model uncertainty.
        \item Parameter uncertainty.
        \item Support class uncertainty.
        \item Model uncertainty.
        \item Construction variability.
    \end{itemize}
    \item Linking SRL East and Snowy 2.0 as recurring design cases.
    \item Producing a defensible design recommendation.
    \item Reporting structure for student and practitioner design reports.
\end{enumerate}

\section{Proposed Figures and Tables}
\begin{itemize}
    \item Integrated workflow diagram.
    \item Method comparison matrix.
    \item Uncertainty register.
    \item Design decision traceability table.
\end{itemize}

\section{Worked Examples and Exercises}
\begin{itemize}
    \item Resolve a case where Q-system support is less conservative than CCM results.
    \item Prepare an integrated design workflow figure for a tunnel project.
    \item Write a short design reflection on uncertainty and additional data needs.
\end{itemize}

\part{Construction Methods, Support Systems and Observational Delivery}

\chapter{Construction Methods, Support Systems and Monitoring}

\section{Chapter Purpose}
Close the design loop by connecting design assumptions to practical excavation methods, support installation, construction sequencing and monitoring-based design adjustment.

\section{Core Content Outline}
\begin{enumerate}
    \item Construction method selection:
    \begin{itemize}
        \item Ground conditions.
        \item Excavation geometry.
        \item Depth and stress environment.
        \item Groundwater.
        \item Settlement tolerance.
        \item Access, program and cost.
        \item Support installation requirements.
    \end{itemize}
    \item Conventional tunnelling:
    \begin{itemize}
        \item Cyclic excavation.
        \item Drill-and-blast.
        \item Mechanical excavation.
        \item Mucking.
        \item Primary support installation.
    \end{itemize}
    \item Mechanised tunnelling:
    \begin{itemize}
        \item TBM principles.
        \item Open hard-rock TBMs.
        \item Shield TBMs.
        \item Earth pressure balance machines.
        \item Slurry shields.
        \item Precast segmental lining.
    \end{itemize}
    \item Surface construction methods:
    \begin{itemize}
        \item Cut-and-cover sequence.
        \item Retaining walls.
        \item Excavation support.
        \item Surface reinstatement.
    \end{itemize}
    \item Immersed tube tunnelling:
    \begin{itemize}
        \item Trench preparation.
        \item Tube immersion.
        \item Jointing and dewatering.
        \item Backfill and protection.
    \end{itemize}
    \item Design and construction philosophies:
    \begin{itemize}
        \item NATM, SEM and SCL.
        \item NMT and Q-system support.
        \item CCM-based ground--support interaction.
        \item ADECO-RS and advance-core control.
    \end{itemize}
    \item Comparing NATM, NMT, CCM and ADECO-RS:
    \begin{itemize}
        \item Ground treated as a load-bearing element.
        \item Single-shell and double-shell support concepts.
        \item Deformation control.
        \item Construction feedback and dynamic design.
    \end{itemize}
    \item Underground support systems:
    \begin{itemize}
        \item Shotcrete.
        \item Fibre-reinforced shotcrete.
        \item Reinforced ribs of sprayed concrete.
        \item Rock bolts and anchors.
        \item Steel ribs and lattice girders.
        \item Forepoling and pipe umbrellas.
        \item Grouting and ground freezing.
        \item Precast concrete segments.
        \item Permanent lining concepts.
    \end{itemize}
    \item Monitoring systems:
    \begin{itemize}
        \item Surface settlement.
        \item Ground movement.
        \item Tunnel convergence.
        \item Ground/support pressure.
        \item Groundwater monitoring.
        \item Trigger levels and response plans.
    \end{itemize}
    \item Construction case integration:
    \begin{itemize}
        \item TBM and drill-and-blast comparison for Snowy 2.0.
        \item Cavern geometry and support demand.
        \item Construction method implications for disturbance factor and support demand.
    \end{itemize}
\end{enumerate}

\section{Proposed Figures and Tables}
\begin{itemize}
    \item Construction method selection matrix.
    \item Conventional versus mechanised tunnelling comparison.
    \item NATM, NMT, CCM and ADECO-RS comparison table.
    \item Temporary and permanent support systems diagram.
    \item Monitoring instrumentation and trigger-action framework.
\end{itemize}

\section{Worked Examples and Exercises}
\begin{itemize}
    \item Select a construction method for shallow urban and deep hard-rock cases.
    \item Explain how the construction method changes disturbance factor and support demand.
    \item Design a monitoring plan linked to expected failure mechanisms and model predictions.
\end{itemize}

\part{Back Matter}

\appendix

\chapter{Calculation Templates}
\begin{itemize}
    \item In-situ stress calculation template.
    \item RMR calculation template.
    \item Q-system calculation template.
    \item GSI comparison template.
    \item Hoek--Brown rock mass parameter template.
    \item Equivalent Mohr--Coulomb fitting template.
    \item CCM calculation template.
    \item Settlement prediction template.
\end{itemize}

\chapter{Worked Project Data Sets}
\begin{itemize}
    \item Melbourne Metro Tunnel context data.
    \item SRL East investigation and modelling data.
    \item Snowy 2.0 underground design data.
    \item Railway station cavern support design data.
\end{itemize}

\chapter{Software and Computational Tools}
\begin{itemize}
    \item Spreadsheet calculations for Hoek--Brown, Mohr--Coulomb and CCM.
    \item RocSupport-type analysis for GRC, SCC and LDP checks.
    \item RS2-type finite element modelling.
    \item UnWedge-type underground block analysis as an optional extension for cavern design.
    \item Python scripts for teaching demonstrations and reproducible calculations.
\end{itemize}

\chapter{Review Questions for the Textbook Outline}
\begin{enumerate}
    \item Should site investigation remain Chapter 2, or should rock mechanics be introduced before ground modelling?
    \item Should empirical support design and failure mechanisms be one chapter, or split into two chapters?
    \item Should Chapter 10 be one large construction chapter, or split into construction methods and observational monitoring?
    \item Should assignment tasks remain as end-of-chapter exercises, or be moved into a separate project workbook?
    \item Should software workflows be included inside the main chapters, or consolidated into appendices?
\end{enumerate}

\chapter{Selected Benchmarking Sources}
\begin{thebibliography}{9}

\bibitem{hoekbrown1980}
E. Hoek and E. T. Brown.
\textit{Underground Excavations in Rock}.
Routledge/Spon Press.
\url{https://www.dymocks.com.au/underground-excavations-in-rock-by-evert-hoek-and-ted-brown-9780419160304}

\bibitem{hudsonharrison1997}
J. A. Hudson and J. P. Harrison.
\textit{Engineering Rock Mechanics: An Introduction to the Principles}.
Elsevier, 1997.
\url{https://shop.elsevier.com/books/engineering-rock-mechanics/hudson/978-0-08-043864-1}

\bibitem{bradybrown2006}
B. H. G. Brady and E. T. Brown.
\textit{Rock Mechanics: For Underground Mining}.
Springer, 2006.
\url{https://link.springer.com/book/10.1007/978-1-4020-2116-9}

\bibitem{kolymbas2005}
D. Kolymbas.
\textit{Tunnelling and Tunnel Mechanics: A Rational Approach to Tunnelling}.
Springer, 2005.
\url{https://link.springer.com/book/10.1007/3-540-28500-8}

\bibitem{kuesel2012}
T. R. Kuesel, E. H. King and J. O. Bickel.
\textit{Tunnel Engineering Handbook}.
Springer, 2012.
\url{https://books.google.com/books/about/Tunnel_Engineering_Handbook.html?id=59V5BgAAQBAJ}

\bibitem{hemphill2012}
G. B. Hemphill.
\textit{Practical Tunnel Construction}.
John Wiley \& Sons, 2012.
\url{https://books.google.com/books/about/Practical_Tunnel_Construction.html?id=exsarm92-64C}

\bibitem{lunardi2008}
P. Lunardi.
\textit{Design and Construction of Tunnels: Analysis of Controlled Deformations in Rock and Soils (ADECO-RS)}.
Springer, 2008.
\url{https://link.springer.com/book/10.1007/978-3-540-73875-6}

\bibitem{rocksoiladeco}
Rocksoil.
\textit{ADECO-RS}.
\url{https://www.rocksoil.com/index.php/en/research-and-development/29-research-and-development/110-01-adeco-rs}

\end{thebibliography}

\end{document}
